\documentclass[11pt]{article}
\usepackage[margin=0.9in]{geometry}
\usepackage{amsmath,amssymb,booktabs,array}
\usepackage[T1]{fontenc}
\usepackage{lmodern}
\usepackage[hidelinks]{hyperref}
\usepackage{microtype}
\setlength{\parindent}{0pt}
\setlength{\parskip}{5pt}
\setlength{\emergencystretch}{2em}
\newcommand{\Lset}{\mathcal L}
\newcommand{\entry}{\mathrm{ent}}
\newcommand{\virt}{\mathrm{virt}}
\newcommand{\wrstate}{\mathrm{wr}}
\newcommand{\can}{\mathrm{can}}
\newcommand{\KL}{\operatorname{KL}}
\newcommand{\tr}{\operatorname{tr}}
\newcommand{\sym}{\operatorname{sym}}
\newcommand{\norm}[1]{\left\lVert#1\right\rVert}
\newcommand{\R}{\mathbb R}
\hypersetup{pdftitle={JLZ v4: Native Joint Local-Delta Optimization and History-Aware Edit Allocation},
 pdfauthor={ODE-edit},pdfsubject={Method specification; not an experimental result}}
\title{JLZ v4: Native Joint Local-$\delta$ Optimization\\
and History-Aware Edit Allocation}
\author{ODE-edit method specification}
\date{October 2, 2026}

\begin{document}
\maketitle

\section{Method claim and design boundary}

JLZ v4 jointly optimizes a local subject intervention at every eligible
write layer. Its methodological claim concerns joint layerwise $z$
optimization and the resulting dynamic allocation of edit effort. The
current-request training sentences remain exactly those of the native
baseline: six native rewrite contexts and one native KL input per request.
There is no additional paraphrase curriculum or generation-prompt training.

The free variables are
\[
 \delta=\{\delta_l\in\R^{d\times B}:l\in\Lset\},
 \qquad\Lset=\{4,5,6,7,8\}.
\]
All five layers are eligible from initialization. Their contributions are
determined jointly through the same nonlinear model loss, rather than
selected in advance. A zero layer or a solution concentrated in one layer
remains permitted. A common subject vector for a given request and layer
is injected across all native contexts.

The $z$ stage uses one model path: a frozen batch-entry model with joint
subject-position injections. The free variables are direct local
activation changes, not writer coordinates $R$. The norm and clamp act on
these injected changes. This design does not use a mixture of an actual
weight branch and an auxiliary branch, an $R$-norm penalty, a realization
consistency loss, or added general/replay prompt losses.

Two arms differ only in an additional history-aware writer value cost.
Arm A ($\eta=0$) is the native joint-injection control. Arm B ($\eta=1$)
adds the value of a regularized local writer fit to the $z$ objective. This
cost includes unrealized target error as well as write energy, so a writer
that rejects a requested change is not incorrectly scored as cheap merely
because its update is small.

After joint optimization, a lower-to-upper writer realizes local targets
at their corresponding layers. The $z$ path and the actual weight model
are distinct. This is an explicit proxy-to-materialization boundary, not
a claim of exact all-token fit/commit equivalence. Joint virtual
interactions are included in the $z$ gradients; actual cross-request and
all-token write effects are only approximated by geometry before writing.

This is a method specification, not a report of implemented GPU fitting,
measured acceleration, or established generalization gains.

\clearpage
\section{Unchanged native inputs and same-site local targets}

Let request $r$ specify a subject, relation prompt, and desired target token
sequence $y_r^*=(y_{r,1}^*,\ldots,y_{r,T_r}^*)$. Preserve native target
normalization, token IDs, teacher-forced prediction positions, context
order, padding masks, and subject-last lookup. The six rewrite inputs are
the canonical prompt and the same five native context wrappers. Their
rewrite losses have equal weights $1/6$.

The native KL input is \texttt{\{\} is a}, with the subject inserted. KL is
measured at its native subject lookup position using the full vocabulary.
Its teacher is frozen at the own batch entry, not at another arm's edited
state. Neither official paraphrase/neighborhood prompts nor raw
\texttt{generation\_prompts} enter the fitting objective.

At layer $l$, write
$v_{lrc}(\delta_{<l})$ for the block output at the subject position
\emph{before the layer's own injection}. It already includes the effects
of all lower-layer interventions in the current joint forward. Define
\begin{equation}
 z_{lrc}(\delta)
 =v_{lrc}(\delta_{<l})+\delta_l[:,r].
 \label{eq:local-target}
\end{equation}
The same $\delta_l[:,r]$ is added at the native block-output hook for every
rewrite context and the KL input of request $r$. Using the same native
hook site also avoids silently changing floating-point addition order by
moving an injection from block output to an earlier linear operation.

All layer injections occur in one forward pass. Higher-layer inputs and
their gradients respond to lower injections. There is no independent
native-$z$ solve per layer, no detach between edited layers, and no fixed
upper-layer activation target used during optimization.

Separately capture the clean entry anchor
\begin{equation}
 h_{lr}^{0}=v_{lr,\can}(0),\qquad a_{lr}^2=\norm{h_{lr}^{0}}_2^2>0.
 \label{eq:anchor}
\end{equation}
The anchor is used only to set the norm scale and clamp radius. In
particular, the moving local target is \emph{not}
$h_{lr}^{0}+\delta_l[:,r]$ once lower interventions are nonzero.
Substituting that expression would remove the required dependence on the
joint lower-layer trajectory.

The local $z$ values and the injected $\delta$ values are distinct objects:
$z$ includes the current local state, while $\delta$ is the directly
optimized displacement. The eventual requested writer residual is exactly
this optimized $\delta_l$ under the payload-preserving contract below.
Its realized displacement can differ because fitting and materialization
are imperfect.

\clearpage
\section{Native joint objective and its single-layer limit}

Let $p_\delta$ denote the output distribution of the jointly injected entry
model. The request loss is
\begin{align}
 n_r(\delta)&=-\frac16\sum_{c=1}^{6}\frac1{T_r}
             \sum_{j=1}^{T_r}
             \log p_\delta(y_{r,j}^*\mid x_{rc},y_{r,<j}^*),\\
 k_r(\delta)&=\KL\!\left(p_{\delta,r,\mathrm{KL}}
                              \middle\|p_{\entry,r,\mathrm{KL}}\right).
\end{align}
Thus the KL direction remains native current relative to entry. The
target-token and context means are followed by a \emph{sum} over requests.
There is one native KL term per request, not one duplicate per layer.

The native joint core is
\begin{equation}
 L_{\mathrm{core}}(\delta)
 =\sum_{r=1}^{B}\bigl[n_r(\delta)+0.0625\,k_r(\delta)\bigr]
   +\sum_{l\in\Lset}\sum_{r=1}^{B}
      \frac{0.5\,\norm{\delta_l[:,r]}_2}{a_{lr}^2},
 \label{eq:core}
\end{equation}
subject to
\begin{equation}
 \norm{\delta_l[:,r]}_2\leq0.75\,a_{lr}.
 \label{eq:clamp}
\end{equation}
Both the denominator and the radius use fixed entry anchors. A candidate
cannot reduce its relative penalty by enlarging a moving normalization
state. There is no norm cost on a subsequently fitted writer residual and
no second regularizer on the same injected vector under another name.

The layer-wise sum extends the native single-vector penalty to the jointly
optimized layer collection. It does not impose equal norms or a minimum
number of nonzero layers. The final model loss determines whether a
request is better served by one layer or several interacting layers.

For a mathematical reduction check, set all $\delta_l=0$ except at one
layer $j$, and set $\eta=0$. With the same inputs, hook site, teacher,
anchor, norm coefficient, clamp, and loss readout, Eq.~\eqref{eq:core}
reduces to the native $z$ objective configured at layer $j$. This is not
necessarily the historical native configuration whose $z$ layer was L8.
It is an objective-level statement, not a claim that a different optimizer,
stopping policy, or subsequent multilayer writer reproduces a native
trajectory exactly. Production eligibility remains all five layers.

In Arm A the losses are separable across requests at the injection stage,
but coupled across layers within a request by the nonlinear forward.
Arm B additionally couples request columns through its small writer
geometry matrices. The actual weight writer shares updates across requests
in either arm; that interference is not fully represented by an
activation-only loss.

\clearpage
\section{History-aware value of a local writer fit}

For each layer, capture target-free entry input keys at the corresponding
subject locations. Preserve native mean-of-group-means pooling:
\begin{align}
 K_l[:,r]&=\frac12 k_{lr1}^{\entry}
                  +\frac1{10}\sum_{c=2}^{6}k_{lrc}^{\entry},\\
 A_l&=15000\,C_{0,l}+H_l^{\entry},\\
 P_l&=(A_l+K_lK_l^\top)^{-1}K_l,\\
 Q_l&=\sym(K_l^\top P_l),\qquad
 s_l^2=\frac1B\sum_r a_{lr}^2.
 \label{eq:geometry}
\end{align}
Here $K_l,P_l\in\R^{m\times B}$, $Q_l\in\R^{B\times B}$,
$C_{0,l}$ is the normalized general-key second moment, and $H_l$ is
post-write key history. The model is FP32 and the geometry is FP64.
The entry geometry is fixed throughout the $z$ solve.

For a desired local displacement matrix $D\in\R^{d\times B}$, define the
regularized writer problem
\begin{equation}
 V_l(D)=\min_{U\in\R^{d\times m}}
 \left\{\frac12\norm{UK_l-D}_F^2
             +\frac12\tr(UA_lU^\top)\right\}.
 \label{eq:value-problem}
\end{equation}
With symmetric positive geometry, stationarity gives
\begin{align}
 U_l^*(D)(A_l+K_lK_l^\top)&=DK_l^\top,\\
 U_l^*(D)&=DP_l^\top,\qquad U_l^*(D)K_l=DQ_l.
\end{align}
The identities
\[
 P_l^\top A_lP_l=Q_l-Q_l^2,\qquad
 (I-Q_l)^2+(Q_l-Q_l^2)=I-Q_l
\]
yield the compact value
\begin{equation}
 \boxed{V_l(D)=\frac12\tr\!\left[D(I-Q_l)D^\top\right]},\qquad
 \nabla_D V_l(D)=D(I-Q_l).
 \label{eq:value}
\end{equation}
The displayed equalities are real-arithmetic identities. Symmetrization in
FP64 fixes the declared numerical operator; small solve/reduction errors
are reported rather than confused with new scientific terms.
Compute $I-Q$, the value, and its analytic gradient in FP64, then cast the
policy gradient to the FP32 intervention dtype for accumulation.

The missed-target component in Eq.~\eqref{eq:value-problem} matters. Along
an eigen-direction with realization $q$, the cost is proportional to
$1-q$. When $q\approx0$, a tiny fitted weight update does not make the
requested displacement cheap: its unrealized target error remains. The
method does not invert $Q_l$ or amplify an ill-conditioned gain to force
perfect realization.

\clearpage
\section{Two allocation arms and the proxy boundary}

Use the injected local changes themselves as the entry writer's targets:
$D=\delta_l$. The complete $z$-stage objective is
\begin{equation}
 \boxed{L_\eta(\delta)=L_{\mathrm{core}}(\delta)
                +\eta\sum_{l\in\Lset}\frac{V_l(\delta_l)}{s_l^2}},
 \qquad \eta\in\{0,1\}.
 \label{eq:arms}
\end{equation}
Its additional smooth gradient is
\begin{equation}
 \nabla_{\delta_l}\left(\eta V_l(\delta_l)/s_l^2\right)
       =\eta\,\delta_l(I-Q_l)/s_l^2.
\end{equation}
The small matrix multiplication uses already prepared geometry and needs
no extra transformer forward or backward.

\begin{center}
\begin{tabular}{p{0.10\linewidth}p{0.30\linewidth}p{0.48\linewidth}}
\toprule
Arm & Objective difference & Interpretation\\
\midrule
A & $\eta=0$ & Native joint local-injection control.\\
B & $\eta=1$ & Adds history-aware local writer value cost.\\
\bottomrule
\end{tabular}
\end{center}

All other inputs, layers, anchors, native coefficients, precision, solver
settings, call budgets, writer rules, and evaluation definitions are
common. The value $\eta=1$ is an initial exploratory coefficient, not a
calibrated performance optimum. Different total objectives cannot be used
as a direct success ranking between the arms.

The additional cost encourages allocation toward changes that the local
writer can realize with less regularized error at the batch entry. It does
not force uniform usage. A single effective layer can receive most of the
change; zero changes elsewhere remain feasible. The group norm and value
cost can also reduce overall edit magnitude, so a smaller cost by itself
is not a successful edit.

The fixed entry $K_l$ is a deliberate approximation. In the eventual
lower-to-upper write, realized lower edits change the upper-layer keys.
The requested local residual remains $\delta_l$; its realized action and
the upper-layer input state can differ from the virtual path. Consequently, $V_l(\delta_l)$
is an allocation proxy, not the exact cost of the final sequential weight
update. It captures local key-space preservation and missed-target error,
not every downstream neighbor or all-token model interaction.

History influences this proxy through $H_l$ and influences the actual
writer through its own current-state solve. There is no additional general
reference KL or replay NLL in Eq.~\eqref{eq:arms}. Geometry-based
preservation must not be described as functional reference-set training.

If $Q_l\ll I$, then $V_l(\delta_l)\approx\norm{\delta_l}_F^2/2$.
In that regime Arm B behaves mainly as additional quadratic regularization;
the presence of $H_l$ in the construction does not establish strong
history-sensitive allocation in the realized run.

\clearpage
\section{One fixed-budget joint Adam solve}

Initialize every $\delta_l$ at zero and optimize all layer/request blocks
with Adam (optimizer weight decay zero): learning rate $0.1$, $\beta_1=0.9$, $\beta_2=0.999$, and
$\epsilon=10^{-8}$. Use 25 joint loss evaluations and 24 updates: evaluate
the initial point, update after each of the first 24 evaluations, and
evaluate the resulting last candidate once. After every update, project
each block onto Eq.~\eqref{eq:clamp}. The PyTorch norm subgradient at zero
is taken as zero, matching the native implementation convention.
Keep zero-valued interventions differentiable: skipping their hooks at
initialization would incorrectly remove their task gradients.

Every logical loss evaluation covers all six rewrite rows plus the KL row
for every request. Microbatch gradients are accumulated before an Adam
step; microbatch size does not alter the loss weights. No per-request
inactive-column mask, per-layer stopping rule, or independently optimized
native-$z$ warm start is used. Both arms have the same fixed schedule.
In particular, this is not a claim of reproducing native per-request
early-stop trajectories, even in a single-layer reduction.

The returned displacement is the last finite clamped candidate. Low PS,
layer concentration, or nonconvergence is recorded rather than used as a
quality gate. Invalid input/state or nonfinite final values are technical
failures. Evaluation prompts never select the returned candidate or
change coefficients during fitting.

\section{Payload-preserving lower-to-upper realization}

Freeze the optimized payloads $\delta_l^*$. Visit L4 through L8. Before
writing layer $l$, run the current actual model, which already contains
the lower-layer writes, to obtain current target-free subject keys.
Pool them with the unchanged native $1/2,1/10$ weights to form
$K_l^{\wrstate}$. With the same batch-entry $A_l$, solve
\begin{align}
 P_l^{\wrstate}&=(A_l+K_l^{\wrstate}(K_l^{\wrstate})^\top)^{-1}K_l^{\wrstate},\\
 U_l^{\wrstate}&=\delta_l^*(P_l^{\wrstate})^\top,\\
 W_l^{\mathrm{new}}&=\operatorname{add}_{32}
             \left(W_l^{\entry},\operatorname{cast}_{32}(U_l^{\wrstate})\right).
\end{align}
The local requested output is the \emph{current} local output plus
$\delta_l^*$. The payload is not changed to compensate for an absolute
virtual target, an L8 residual, or a remaining-layer divisor. There is no
second allocation step and no inverse-$Q$ gain correction.

In ideal arithmetic, the realized displacement on the pooled current keys
is $\delta_l^*Q_l^{\wrstate}$, with
$Q_l^{\wrstate}=(K_l^{\wrstate})^\top P_l^{\wrstate}$, rather than necessarily
$\delta_l^*$. Moreover $Q_l^{\wrstate}$ need not equal the entry proxy $Q_l$.
Finite-precision materialization, individual contexts, non-subject token
effects, and lower-layer realization errors create further differences
from the virtual joint path. The specification does not hide these gaps
inside a claim of exact realization.
Measure materialized action using $(W_l^{\mathrm{new}}-W_l^{\entry})K_l^{\wrstate}$
with the subtraction evaluated in FP64. This differs from using only the
cast update; actual linear-output subtraction can additionally reflect
GEMM rounding.

After all five writes, recapture post-write keys on the fully edited model
and append each $H_l$ once. $H_l$ is fixed during all entry and current-key
solves for that batch. Each arm follows its own committed sequential state.
A zero payload writes a zero update at its layer even if lower writes have
changed that layer's inputs; no compensating edit is introduced there.

\clearpage
\section{Compute intent, limitations, and reporting}

The $z$ stage has one joint subject-injection model path: at BS100 each
evaluation has 700 native input rows, with 24 forward/backward evaluations
and one final forward evaluation. There is no second full-weight branch,
no extra general/replay prompt batch, and no large weight-gradient tensor
for the frozen model. The small value-cost gradient is added analytically.
All five $\delta_l$ remain jointly optimized; no layer subset saves compute.

This removes the prior two-path fitting structure, but is not a measured
speed claim. Arm B prepares entry geometry once for the proxy, and the final
writer additionally fits the current keys once per layer. Key capture,
linear solves, the model's input-length distribution, Adam state, final
observations, and history updates must be included in total time and memory.
An entry solve may be reused only when its current key and geometry
identities actually match, not merely because both solves have the same shape.

The dynamic allocation is driven by joint virtual losses, per-layer
native regularization, and, in Arm B, approximate writer value. A small
norm or a concentrated allocation is neither a success nor a failure by
itself. Compare actual committed RS/PS/NS and sequential retention with
fixed denominators; distinguish virtual fit from actual performance.
Report norm, clamp use, proxy cost, current writer cost, and compute
separately. Different hardware or stopping schedules preclude a matched
runtime claim from historical baseline numbers alone.

Local alignment removes the need to reinterpret every layer's target as
an L8 vector, but it does not ensure that a closed-form write exactly
realizes the optimized joint path. Shared subject injections also do not
guarantee unseen-paraphrase generalization. The entry proxy may become
inaccurate after lower writes, and $\eta=1$ may suppress useful strength.
These are method limits, not reasons to add strict quality or layer gates.

This document specifies a proposed native joint-local-$\delta$ method.
It does not modify an existing execution lock or provide a v4 runner.
Implementation, GPU validation, and performance measurements are not
claimed here.

\textbf{Primary-source context.}
\href{https://papers.neurips.cc/paper_files/paper/2025/file/70d4ef44dc973586cfa3ea92b4868b72-Paper-Conference.pdf}{BLUE}
uses boundary-layer targets;
\href{https://aclanthology.org/2026.acl-long.918.pdf}{CAKE}
uses causal-weighted residual assignment;
\href{https://arxiv.org/html/2605.00358v2}{FE}
constructs compatible targets by forward replay from a first-layer anchor;
\href{https://aclanthology.org/2026.acl-long.1855.pdf}{HiEdit}
combines lifelong preservation with adaptive sparse layer selection.
Dynamic allocation or preservation alone is not the novelty claim here.

\end{document}
